% !PRINTED 452
\weberpage{452}
The whole system is carried over into itself by 168 substitutions of the group.

\[
\begin{aligned}
&\text{On every line } A \text{ lie four points } a. \\
&\text{Through every point } a \text{ pass four lines } A. \\
&\text{On every line } T \text{ lie three points } a. \\
&\text{Through every point } t \text{ pass three lines } A.
\end{aligned}
\]

\gesperrt{The points} $a$ \gesperrt{and} $t$ \gesperrt{together form the complete system of all points of intersection of the lines} $A$\gesperrt{.}

For 21 straight lines intersect in 210 points, and of these, three or six coincide when three or four of these lines pass through one point. But we have
\[
210 = 21.6 + 28.3.
\]

\gesperrt{On every line} $A$ \gesperrt{lie four points} $t$\gesperrt{.}

For since through every point $t$ there pass three lines $A$, and on every line $A$ equally many points $t$ must lie (because every line $A$ is transformable into every other), it follows, if $x$ is the number of these points, that $x . 21 : 3 = 28$, hence $x = 4$. In the same way one concludes:

\gesperrt{Through every point} $a$ \gesperrt{pass four lines} $T$\gesperrt{.}

None of the points $t$ lies on a line $T$.
If we denote the $\nu$-fold poles by $P_\nu$, we thus have the following systems of poles:
\[
\begin{aligned}
&21 \text{ eightfold poles } P_8, \\
&28 \text{ sixfold} \;\text{\textquotedbl}\; P_6, \\
&42 \text{ fourfold} \;\text{\textquotedbl}\; P_4, \\
&56 \text{ threefold} \;\text{\textquotedbl}\; P_3, \\
&24 \text{ sevenfold} \;\text{\textquotedbl}\; P_7, \\
&\text{infinitely many twofold poles } P_2.
\end{aligned}
\]

The $P_2$ are all the points of the principal axes distinct from $P_4, P_6, P_8$; of the points $P_4$, two at a time lie on a line $A$; of the points $P_3$, two at a time lie on a line $T$; the points $P_7$ do not lie on the lines $A$ (that they also do not lie on $T$ will appear later).

The system of the axes corresponds exactly to the system of the poles. Every straight line that passes through an eightfold pole is an axis (among them the principal axes, four lines $T$ and two connecting lines of the $P_8$ with $P_4$). Accordingly
%CONT>>
